\paragraph{Definitional status of the prescription.} The formulation above takes \(\mathbf r\) and \(H(\mathbf r)\) as data: at fixed \(\mathbf r\) it presupposes a microscopic Hamiltonian depending on separation as an external parameter, so that \(\delta\hat F_i=-\partial H/\partial r_i\) and the correlator are defined. On the bare data of Presentation~A this is not available. The coordinates of Definition~\ref{def:matrix-coords-ext} are matrix-valued and the action Equation~\eqref{eq:Smatrix} contains no pair-separation parameter with respect to which a derivative could be taken, and separation is among the infrared quantities the programme takes to be emergent, so the prescription is not a statement about the matrix model until one of three exits is taken. (i)~A fixed-background class with commuting base coordinates: in the conjugate-coordinate class of the Problem~F entry the base functions \(y^\mu\) commute, so a separation between two localized configurations can be written as a difference in \(y\)-space; but that class must then be extended by a sector supplying localized cores (its sector list contains none), and the identification of a \(y\)-difference with the separation \(d\) and sizes \(\ell_A,\ell_B\) used above is an additional input, so the resulting \(\Gamma_{ij}\) is a coefficient of the extended fixed-background system, not of the bare matrix model. (ii)~The Presentation~B dictionary (Assumption~\ref{ass:dictionary}): geometric separations are available there, but every calculation inherits the dictionary, and a \(\Gamma_{ij}\) obtained this way carries the dictionary's conditionality and cannot discharge the matrix-derived label. (iii)~An explicitly reduced model with imported separation: a bath Hamiltonian, coupling and the parameter \(\mathbf r\) are supplied as inputs; the prescription is then well posed, but the model is an assumption and its \(\Gamma_{ij}\) constrains the NAMCG correlator only if the dictionary route of (ii) is also taken. These statements are for the bare Presentation-A data and the three exits as characterized here; they compute no response, do not exclude the existence of a further route, and change no ledger label; Problem~I remains open.

\paragraph{Low-frequency support and sub-critical vanishing (conditional; not closure).} Within exit~(iii), where the data are supplied, the prescription itself restricts the answer. Let the reference state be stationary and passive, diagonal in the eigenbasis of \(H\) with weights \(p_n\), and let \(F^i_{mn}=\langle m|\delta\hat F_i|n\rangle\) at fixed \(\mathbf r\). With the transform convention of the display above, \(\chi^R_{ij}(\omega)=\int\mathrm dt\,e^{i\omega t}\chi^R_{ij}(t)\), the transform of the commutator has a principal-value real part and, for \(\omega>0\), an imaginary part
\[
\operatorname{Im}\chi^R_{(ij)}(\omega;\mathbf r)
=
\frac{\pi}{\hbar}
\sum_{m,n}
(p_n-p_m)\,
\operatorname{Re}\!\bigl(F^i_{mn}F^j_{nm}\bigr)\,
\delta\!\left(\omega-\frac{E_m-E_n}{\hbar}\right),
\]
so \(\operatorname{Im}\chi^R_{(ij)}\) is supported on the weighted Bohr frequencies of the connected force operators. Consequently, whenever that support stays a positive distance from \(\omega=0\) -- a gapped reference, or finite volume -- \(\Gamma_{ij}=0\) exactly; and for a gapless bath in the thermodynamic limit the tensor vanishes when the weighted density of states vanishes faster than \(\omega\) at low frequency, e.g.\ for finite-size Gaussian cores of size \(\ell\) with derivative coupling to a linear bath \(\varepsilon_q=cq\) in three spatial dimensions, where the form-factor integral gives \(\operatorname{Im}\chi^R(\omega)\propto\omega^3 e^{-\omega^2\ell^2/c^2}\), and for gapped baths \(\varepsilon_q=\sqrt{\Delta^2+c^2q^2}\), where the support starts at \(\omega=\Delta\). The vanishing is not universal: with derivative coupling in one spatial dimension the same integral gives \(\operatorname{Im}\chi^R(\omega)\propto\omega e^{-\omega^2\ell^2/c^2}\) and \(\Gamma\neq0\), so the question is exactly the low-frequency weighted spectral weight. Distinct from this \(\omega\downarrow0\) coefficient is the on-shell absorption rate at velocity \(v\), supported on \(\omega=v\,\mathbf q\cdot\hat{\mathbf v}\): it vanishes for \(v<\min_q\varepsilon_q/|\mathbf q|\) -- on a lattice the minimum phase velocity, not the long-wavelength sound speed, sets the threshold -- which is the Landau criterion \cite{Landau1941}, a kinematic criterion and not a theorem; a medium-modified dispersion with lowered low-momentum phase velocity reopens sub-critical Cherenkov friction \cite{BerezhianiElderKhoury2019}, the same mechanism family as modified propagation in the Problem~F entry. A drag law of the working form \(\Gamma\propto d^{-4}\) with nonzero coefficient (Assumption~\ref{ass:drag}) therefore requires, beyond a definable separation, an open low-frequency channel with nonzero weighted spectral weight at \(\omega\downarrow0\) at every separation in the regime, and its exponent and tensor structure are fixed by that channel, not by any static tail. These statements are for the exit-(iii) class of supplied Hamiltonians, stationary passive reference states and finite-size Gaussian cores at the stated truncations; they do not compute the NAMCG correlator, do not constrain the dictionary or the emergent metric, do not establish or exclude the \(d^{-4}\) law, and give no value or sign of any coefficient; Problem~I remains open.
